\chapter{2007年上海卷（秋文）}
\section{填空题}
\begin{problem}
方程 \( 3^{x-1} = \frac{1}{9} \) 的解是 \fillinblank{} \fillinblank{}.
\end{problem}

\begin{problem}
函数 \( f(x)=\frac{1}{\pi-1} \) 的反函数 \( f^{-1}(x)= \)  \fillinblank{}.
\end{problem}

\begin{problem}
直线 \(4x + y - 1 = 0\) 的倾斜角 \(\theta =\)  \fillinblank{}.
\end{problem}

\begin{problem}
函数 \( y = \sec x \cdot \cos \left( x + \frac{\pi}{4} \right) \) 的最小正周期 T =\fillinblank{} \fillinblank{}.
\end{problem}

\begin{problem}
以双曲线 \(\frac{x^2}{4} - \frac{y^2}{5} = 1\) 的中心为顶点，且以该双曲线的右焦点为焦点的抛物线方程是 \fillinblank{} \fillinblank{}.
\end{problem}

\begin{problem}
若向量 \(\bs{a}\)，\(\bs{b}\) 的夹角为 \(60^{\circ}\)，\(|\bs{a}| = |\bs{b}| = 1\)，则 \(\bs{a} \cdot (\bs{a} - \bs{b}) = \)  \fillinblank{}.
\end{problem}

\begin{problem}
如图，在直三棱柱 \(ABC - A_{1}B_{1}C_{1}\) 中，\(\angle ACB = 90^{\circ}\)，\(AA_{1} = 2\)，\(AC = BC = 1\)，则异面直线 \(A_{1}B\) 与 \(AC\) 所成角的大小是 \fillinblank{}. (结果用反三角 函数值表示)
\end{problem}

\begin{problem}
某工程由 A, B, C, D 四道工序组成，完成它们需用时间依次为 2, 5, x, 4 天. 四道工序的先后顺序及相互关系是: A, B 可以同时开工; A 完成后，C 可以开工; BC 完成后，D 可以开工. 若该工程总时数为 9 天，则完成工序 C 需要的天数 x 最大是 \fillinblank{} \fillinblank{}.
\end{problem}

\begin{problem}
在五个数字 1, 2, 3, 4, 5 中，若随机取出三个数字，则剩下两个数字都是奇数的概率是 \fillinblank{}. (结果用数值表示)
\end{problem}

\begin{problem}
对于非零实数 \(a, b\)，以下四个命题都成立: \circled{1} \(a + \frac{1}{a} \neq 0\); \circled{2} \((a + b)^2 = a^2 + 2ab + b^2\); \circled{3} 若 \(|a| = |b|\)，则 \(a = \pm b\); \circled{4} 若 \(a^2 = ab\)，则 \(a = b\). 那么，对于非零复数 \(a, b\)，仍然成立的命题的所有序号是 \fillinblank{} \fillinblank{}.
\end{problem}

\begin{problem}
如图，\(AB\) 是直线 \(l\) 上的两点，且 \(AB = 2\). 两个半径相等的动圆分别与 \(l\) 相切于 \(AB\) 点，\(C\) 是这两个圆的公共点，则圆弧 \(AC, CB\) 与线段 \(AB\) 围成图形面积 \(S\) 的取值范围是 \fillinblank{} \fillinblank{}.
\end{problem}

\section{单选题}
\begin{problem}
已知 a, b \(\in\) R，且 \( 2 + a, b + 3\i \) (i 是虚数单位) 是一个实系数一元二次方程的两个根，那么 a, b 的值分别是（\quad）

\choices
    {\(a = -3, b = 2\)}
    {\(a = 3, b = -2\)}
    {\(a = -3, b = -2\)}
    {\(a = 3, b = 2\)}
\end{problem}

\begin{problem}
圆 \(x^{2} + y^{2} - 2x - 1 = 0\) 关于直线 \(2x - y + 3 = 0\) 对称的圆的方程是（\quad）

\choices
    {\((x + 3)^2 + (y - 2)^2 = \frac{1}{2}\)}
    {\((x - 3)^2 + (y + 2)^2 = \frac{1}{2}\)}
    {\((x + 3)^2 +(y - 2)^2 = 2\)}
    {\((x - 3)^2 +(y + 2)^2 = 2\)}
\end{problem}

\begin{problem}
数列 \(\{a_{n}\}\) 中，\(a_{n} = \begin{cases} \frac{1}{n^{2}}, & 1 \leqslant n \leqslant 1000, \\ \frac{n^{2}}{n^{2} - 2n}, & n \geqslant 1001, \end{cases}\) 则数列 \(\{a_{n}\}\) 的极限值（\quad）

\choices
    {等于 0}
    {等于 1}
    {等于 0 或 1}
    {不存在}
\end{problem}

\begin{problem}
设 \( f(x) \) 是定义在正整数集上的函数，且 \( f(x) \) 满足：“当 \( f(k) \geq k^2 \) 成立时，总可推出 \( f(k+1) \geq (k+1)^2 \) 成立”. 那么，下列命题总成立的是（\quad）

\choices
    {若 \( f(1) < 1 \) 成立，则 \( f(10) < 100 \) 成立}
    {若 \( f(2) < 4 \) 成立，则 \( f(1) \geqslant 1 \) 成立}
    {若 \( f(3) \geqslant 9 \) 成立，则当 \( k \geqslant 1 \) 时，均有 \( f(k) \geqslant k^2 \) 成立}
    {若 \( f(4) \geqslant 25 \) 成立，则当 \( k \geqslant 4 \) 时，均有 \( f(k) \geqslant k^2 \) 成立}
\end{problem}

\section{解答题}
\begin{problem}
在正四棱锥 \(P - ABCD\) 中，\(PA = 2\)，直线 \(PA\) 与平面 \(ABCD\) 所成的角为 \(60^{\circ}\)，求正四棱锥 \(P - ABCD\) 的体积 \(V\).
\end{problem}

\begin{problem}
在 \(\triangle ABC\) 中，\(a, b, c\) 分别是三个内角 \(A, B, C\) 的对边. 若 \(a = 2, C = \frac{\pi}{4}\)，\(\cos \frac{B}{2} = \frac{2\sqrt{5}}{5}\)，求 \(\triangle ABC\) 的面积 \(S\).
\end{problem}

\begin{problem}
近年来，太阳能技术应用的步伐日益加快. 2002 年全球太阳电池年的年产生量达到 670 兆瓦，年生产量的增长率为 34 点. 以召回年中，年生产量的增长率逐年递增 2 点 (如，2003 年的年生产量的增长率为 36 点). (1) 求 2006 年全球太阳电池的年生产量 (结果精确到 0.1 兆瓦); (2) 目前太阳电池产业存在的主要问题是将市场安装量远小于生产量，2006年的实际安装量为 1420 兆瓦. 假设以后若干年内太阳电池的年生产量的增长率保持在 42\%，到 2010 年，要使年安装量与年生产量基本持平 (即年安装量不少于年生产量的 95\%)，这四年中太阳电池的年安装量的平均增长率至少应达到多少 (结果精确到 0.1\% ？
\end{problem}

\begin{problem}
已知函数 \( f(x) = x^2 + \frac{a}{x} (x \neq 0, \text{常数 } a \in \R) \). (1) 当 \(a = 2\) 时，解不等式 \(f(x) - f(x - 1) > 2x - 1\); (2) 讨论函数 \( f(x) \) 的奇偶性，并说明理由.
\end{problem}

\begin{problem}
如果有穷数列 \(a_1, a_2, a_3, \cdots, a_m\) (m 为正整数) 满足条件 \(a_1 = a_m\)，\(a_2 = a_{m-1}, \cdots, a_m = a_1\)，即 \(a_i = a_{m-i+1}\) (\(i = 1, 2, \cdots, m\))，我们称其为“对称数列”. 例如，数列 1, 2, 5, 2, 1 与数列 8, 4, 2, 2, 4, 8 都是“对称数列”. (1) 设 \(\{b_n\}\) 是项数为 7 的“对称数列”，其中 \(b_1, b_2, b_3, b_4\) 是等差数列，且 \(b_1 = 2, b_4 = 11\). 依次写出 \(\{b_n\}\) 的每一项: (2) 设 \(\{c_{n}\}\) 是 49 项的“对称数列”，其中 \(c_{25}, c_{26}, \cdots, c_{40}\) 是首项为 1，公比为 2 的等比数列，求 \(\{c_{n}\}\) 各项的和 \(S\)； (3) 设 \(\{d_n\}\) 是 100 项的“对称数列”，其中 \(d_{51}, d_{52}, \cdots, d_{100}\) 是首项为 2，公差为 3 的等差数列. 求 \(\{d_n\}\) 前 \(n\) 项的和 \(S_n (n = 1, 2, \cdots, 100)\).
\end{problem}

\begin{problem}
我们把由半椭圆 \(\frac{x^2}{a^2} + \frac{y^2}{b^2} = 1 (x \geqslant 0)\) 与半椭圆 \(\frac{y^2}{b^2} + \frac{x^2}{a^2} = 1 (x \in 0)\) 合成的曲线称作“果圆”，其中 \(a^2 = b^2 + c^2, a > 0, b > c > 0\). 如图，设点 \(F_0, F_1, F_2\) 是相应椭圆的焦点，\(A_1, A_2\) 和 \(B_1, B_2\) 分别是“果圆”与 \(x, y\) 轴的交点，\(M\) 是线段 \(A_1A_2\) 的中点. (1) 若 \(\triangle F_0F_1F_2\) 是边长为 1 的等边三角形，求该“果圆”的方程; (2) 设 \(P\) 是“果圆”的半椭圆 \(\frac{x^2}{B^2} + \frac{y^2}{A^2} = 1 (x \leqslant 0)\) 上任意一点. 求证: 当 \(|PM|\) 取得最小值时，\(P\) 在点 \(B_1, B_2\) 或 \(A_1\) 处; (3) 若 \(P\) 是“果圆”上任意一点，求 \(|PM|\) 取得最小值时点 \(P\) 的横坐标.
\end{problem}

